\documentclass[12pt]{article}
\usepackage{graphicx} % Required for inserting images
\usepackage[margin=0.5in]{geometry}
\usepackage{amsmath, amssymb}
\usepackage{mathrsfs}


\usepackage{soul}


% Dr. David Brown Commands
\newcommand{\ds}{\displaystyle}
\newcommand{\R}{\mathbb{R}}
\newcommand{\N}{\mathbb{N}}
\newcommand{\Q}{\mathbb{Q}}
\newcommand{\C}{\mathbb{C}}


% Commands to get script r
\usepackage{calligra}
\DeclareMathAlphabet{\mathcalligra}{T1}{calligra}{m}{n}
\DeclareFontShape{T1}{calligra}{m}{n}{<->s*[2.2]callig15}{}
\newcommand{\scripty}[1]{\ensuremath{\mathcalligra{#1}}}

% Unit commands
\newcommand{\degree}{\text{°}}
\newcommand{\unitSpeed}{\frac{\text{m}}{\text{s}}}
\newcommand{\unitAcceleration}{\frac{\text{m}}{\text{s}^2}}

% Constant commands
\newcommand{\avogadro}{6.022\times10^{23}}

\newcommand{\boltzmannConstK}{1.381\times10^{-23}\frac{\text{J}}{\text{K}}}
\newcommand{\planckConst}{6.626\times10^{-34}\text{J}\cdot\text{s}}

\newcommand{\atmP}{1.01325\times10^5\,\text{Pa}}

% Commands to easily write partial derivatives
\newcommand{\partDeriv}[2]{\frac{\partial #1}{\partial #2}}
\newcommand{\doublePartDeriv}[2]{\frac{\partial^2 #1}{\partial^2 #2}}


\title{Problem 6.6}
\author{Gabriel Decker}


\begin{document}



\maketitle
\begin{flushleft}



Recall the following equations.\\~\\

First, the equation for the probability for a particular state is
\begin{equation}
    \mathcal{P}(s)=\frac{1}{Z}e^{-E(s)/kT}
\end{equation}
where $Z$ is the following.
\begin{equation}
    Z=\sum_se^{-E(s)/kT}
\end{equation}\\~\\

So first, we'll have to find $Z$ before we can find any probabilities.
The first couple of energy levels are the following for a Hydrogen atom:
$E_1=-13.6\,\text{eV}$, $E_2=-3.4\,\text{eV}$, $E_3=-1.5\,\text{eV}$.
The ground state $n=1$ has multiplicity of 1. The $n=2$ state has multiplicity of 4. The $n=3$ state has multiplicity of 9.\\~\\

Before plugging into the equation for $Z$, let's find $k$ in $\text{eV/K}$.
$$k=\boltzmannConstK\cdot\frac{6.242\times10^{18}\,\text{eV}}{1\,\text{J}}=8.6202\times10^{-5}$$\\~\\

We can start calculating $Z$ by using the first three energy levels.
$$Z\approx e^{13.6\,\text{eV}/(8.6202\times10^{-5}\cdot T)} + 4e^{3.4\,\text{eV}/(8.6202\times10^{-5}\cdot T)} + 9e^{1.5\,\text{eV}/(8.6202\times10^{-5}\cdot T)}$$
For room temp $T=300\,\text{K}$, this becomes
$$\implies Z\approx e^{525.897}+4e^{131.474}+9e^{58.003}$$
Clearly, the first term (the ground state's term) dominates $Z$, so approximate $Z$ as
$$\implies Z\approx e^{-E_1/kT}$$\\~\\

Now that we have an approximation for the partition function $Z$, we can plug it into the equation for the probability that the atom is in the first excited state ($E_2=-3.4\,\text{eV}$). Recall that this state as multiplicity 4.
$$\mathcal{P}(E_2)=4\frac{e^{-E_2/kT}}{e^{-E_1/kT}}$$
$$\implies\mathcal{P}(E_2)=4e^{-(E_2-E_1)/kT}$$
$$\implies\mathcal{P}(E_2)=4e^{-(10.2\,\text{eV})/kT}$$\\~\\

For $T=300\,\text{K}$,
$$\implies\mathcal{P}(E_2)=4e^{-(10.2)/(8.6202\times10^{-5}\cdot 300)}$$
$$\implies\mathcal{P}(E_2)=4e^{-394.4}$$
$$\implies\mathcal{P}(E_2)\approx 0$$
So there is essentially zero chance of finding a Hydrogen atom in its first excited state at room temperature.\\~\\

For $T=9500\,\text{K}$,
$$\implies\mathcal{P}(E_2)=4e^{-(10.2)/(8.6202\times10^{-5}\cdot 9500)}$$
$$\implies\mathcal{P}(E_2)=4e^{-12.46}$$
$$\implies\mathcal{P}(E_2)=0.0000156$$
So the chance of finding a Hydrogen atom in its first excited state is much higher at the surface of $\gamma$ UMa.\\~\\


\end{flushleft}

\end{document}
